\subsection*{Opposed evidence changes the written position}
Can the writer answer the same question differently when the evidence changes? We selected natural Chinese questions with affirmative and negative human responses supported by different paragraphs. Both writers learned to reproduce the same human quotations. Only one received the stance appropriate to each passage throughout training. On 18 official development questions the matched writer gave the correct stance on both passages for 10 questions, versus one for reassigned supervision and six for a fixed lexical marker rule (Fig.~\ref{fig:paired-stance-writer}). On 144 validation questions the counts were 55, 13 and 41. The paired validation advantage over the marker was 9.7 percentage points (95\% bootstrap interval $-0.7$ to 19.4), missing the prespecified superiority criterion. All 144 validation source records had been available in the inherited answer corpus, whereas the 18 official records had not.

Eleven of 36 official responses reproduced the position, quotation and stop, compared with seven under reassignment. Nine of the eleven successes quoted the entire supplied window. An exact human quotation following either valid position occurred 12 times for the matched writer and 13 for its control. The matched writer gave a correct stance on 26 responses requiring at least one stance glyph absent from the input. Its source of evidence was largely the passage itself. Erasing the question left both stances correct for 11 of 18 questions, against 10 with the question present. On a separate constructed test, changing the question reversed the warranted stance on each of 24 passages. The matched writer kept the same stance all 24 times. It attained 24 of 48 individual stances and zero of 12 complete four-way contrasts. These findings identify a passage-dependent writing skill and a persistent failure to bind the question to the relevant claim.
